\documentclass[9pt]{beamer}
\usefonttheme[onlymath]{serif}
% Theme settings
\usetheme{metropolis} % Professional, clean theme
\usepackage[english]{babel}
\usepackage{booktabs}
\usepackage{amsmath}
\usepackage{amssymb}
\usepackage{bm}
\usepackage{mdframed}

% Custom Commands from your report
\newcommand{\ten}[1]{\bm{\mathcal{#1}}}
\newcommand{\mat}[1]{\mathbf{#1}}
\newcommand{\norm}[1]{\left\lVert#1\right\rVert}

\title{Improved Recipe for Compress-then-Train}
\subtitle{Reasoning-aware, training-aware compression with on-policy distillation}
\author{Yequan Zhao}
\date{\today}

\begin{document}
	
	\maketitle
	
	% ---------------------------------------------------------------------------
	\begin{frame}{Proposed intern project}
		
		\textbf{Goal:} a better recipe for \emph{compress-then-train}: optimize for the model \emph{after} training, not only at compression time.
		
		\begin{itemize}
			\item \textbf{Reasoning-aware compression} --- calibrate the compression on full reasoning traces so the model's generation / reasoning ability survives the cut.
			\item \textbf{Training-aware compression} --- pick the compression that
			\emph{converges fastest} once training resumes, not just the one with the
			lowest one-shot error.
			\item \textbf{Extended training recipe} --- combine off-policy knowledge
			distillation with \textbf{on-policy distillation (OPD)} to keep improving
			on reasoning tasks after standard KD plateaus.
		\end{itemize}
		
	\end{frame}

	% ---------------------------------------------------------------------------
	\begin{frame}{Reasoning-aware compression: calibration preserves generation}

		\textbf{Challenge.} One-shot structured low-rank compression (SVD-attn +
		Nystr\"om-MLP) is very weak at preserving \emph{reasoning} --- the
		compressed model loses its ability to terminate a reasoning trace (loops
		past the answer), collapsing math accuracy.

		\vspace{0.4em}
		\textbf{Cause.} The calibration covariances were collected on reasoning
		traces \emph{truncated into fixed-length windows} --- this drops the
		long-trace tails and lets long traces dominate the token pool, so the
		covariance no longer reflects the reasoning distribution.

		\vspace{0.4em}
		\textbf{Fix.} Calibrate on the \emph{full reasoning trace} (no truncation)
		and \emph{reweight by sequence} (every conversation counts equally, not
		every token). Pure calibration change --- same method, same budget.

		\vspace{0.5em}
		{\footnotesize Qwen3-4B$\to$, retain $0.7$, forward-only SVD+Nystr\"om;
		MATH-500 strict (greedy), C4 PPL.}

		\vspace{0.2em}
		{\footnotesize
		\begin{tabular}{@{}lcc@{}}
			\toprule
			calibration scheme & MATH-500 & C4 PPL \\
			\midrule
			truncated windows (2048)            & 66.0\%          & 103.0 \\
			full sequence                       & 66.0\%          & 103.0 \\
			full sequence + sequence-reweight   & \textbf{69.0\%} & \textbf{98.7} \\
			\bottomrule
		\end{tabular}}

		\vspace{0.5em}
		Sequence-reweighting is the decisive axis ($+3$pp MATH, lower PPL); it also
		removes the looping/termination failure (the model reaches the answer
		\emph{and} stops). Now the repo default.

	\end{frame}

	% ---------------------------------------------------------------------------
	\begin{frame}{Training-aware compression}

		Input-output whitened SVD
		\[
		\hat{\mat{W}}
		\;=\;
		\arg\min_{\hat{\mat{W}}}\;
		\norm{\,\mat{C}_g^{1/2}\,(\mat{W}-\hat{\mat{W}})\,\mat{C}_x^{1/2}\,}_F^2 ,
		\]
		solved by truncated SVD of $\mat{M}=\mat{C}_g^{1/2}\,\mat{W}\,\mat{C}_x^{1/2}$,
		with $\mat{C}_x=\sum_t\mat{x}_t\mat{x}_t^{\!\top}$ (forward activations) and
		$\mat{C}_g=\sum_t\mat{g}_t\mat{g}_t^{\!\top}$ (back-propagated loss gradients).
		
		\vspace{0.8em}
		This two-sided whitening is precisely a Hessian estimate for the layer weights:
		\[
		\mat{H}_{\mat{W}} \;=\; \mathbb{E}\!\left[\,\mat{x}\mat{x}^{\!\top}
		\,\otimes\, \mat{g}_y\mat{g}_y^{\!\top}\,\right].
		\]
		
		\vspace{0.4em}
		So truncation discards the \emph{least loss-sensitive} subspace, and those
		surviving high-curvature directions are exactly the ones training keeps
		actively updating.
		
		Nystr\"om: build a single joint forward+backward PSD kernel
		$\mat{K}_{\text{joint}}=\mat{C}_f^{1/2}\mat{C}_b\mat{C}_f^{1/2}$ and select the
		hidden neurons by Nystr\"om approximation: keeping the directions important to
		\emph{both} the forward pass and the gradient that trains $\mat{W}_u,\mat{W}_g$.
		
	\end{frame}
	
	% ---------------------------------------------------------------------------
	\begin{frame}{MoE expert compression: doubly-whitened Nystr\"om}
		

		\vspace{0.5em}
		{\footnotesize OLMoE-1B-7B, retain $0.5$; eval limit 200/task. After-training =
			short SFT recovery ($\sim$10k OpenThoughts samples).}
		
		\vspace{0.2em}
		{\footnotesize
			\begin{tabular}{@{}lcccc@{}}
				\toprule
				& \multicolumn{2}{c}{training-free} & \multicolumn{2}{c}{after training} \\
				\cmidrule(lr){2-3}\cmidrule(lr){4-5}
				method & MMLU & GSM8K & MMLU & GSM8K \\
				\midrule
				original (uncompressed)     & 0.542 & ---   & ---   & ---   \\
				\midrule
				Nystr\"om-combined (fwd+bwd) & \textbf{0.382} & \textbf{0.150} & \emph{TBD} & \emph{TBD} \\
				Nystr\"om (fwd-only)         & 0.341 & 0.135 & \emph{TBD} & \emph{TBD} \\
				REAP (expert drop)           & 0.289 & 0.000 & \emph{TBD} & \emph{TBD} \\
				SlimQwen (expert merge)      & 0.269 & 0.005 & \emph{TBD} & \emph{TBD} \\
				\bottomrule
		\end{tabular}}
		
		\vspace{0.5em}
		\begin{itemize}
			\item \textbf{fwd+bwd $>$ fwd-only $>$ expert removal/merge} training-free:
			joint kernel beats forward-only; both low-rank methods beat coarse
			whole-expert drop/merge (which collapse at retain $0.5$).
			\item \textbf{Open question:} does short training reorder these families?
			Recovery leg is the next run.
		\end{itemize}
		
	\end{frame}
	
	% ---------------------------------------------------------------------------
	\begin{frame}{Training recipe: off-policy KD vs.\ on-policy distillation}
		
		\begin{itemize}
			\item \textbf{Off-policy KD} --- student matches the teacher's traces. 
			\item \textbf{On-policy distillation (OPD)} --- student \emph{generates} the
			responses, teacher \emph{scores} them token-by-token. The student is
			corrected exactly where \emph{it} goes wrong.
		\end{itemize}
		
		\vspace{1em}
		\begin{mdframed}
			\textbf{Two regimes, in order:}
			\begin{itemize}
				\item \textbf{Stage 1 --- off-policy KD is more efficient.} Cheap dense
				supervision recovers most of the quality fast.
				\item \textbf{Stage 2 --- when Stage 1 plateaus, OPD keeps improving} on
				reasoning tasks
				% fixing the student's own failure modes that the teacher's traces never expose.
			\end{itemize}
		\end{mdframed}
		
	\end{frame}
	
	% ---------------------------------------------------------------------------
	\begin{frame}{Broader ideas}
		
		% \textbf{Heterogeneous / adaptive compression ratios.} Allocate the budget \emph{across layers and experts} according to training-aware curvature --- spend rank where it converges best after training, not uniformly.
		
		Compression-aware training:
		\begin{itemize}
			\item Train a 8B model, but contains a good 4B (or more) models -> once for all
			
			\item Full-rank training-aware transformation
			\begin{itemize}
				\item The parameter trainability (curvature) is ranked: BlockTT (rank by singular value) / Nystrom (rank by row/col leverage score)
				\item Regularized training: only the information beyond the 4B model capacity goes to 8B only parameters.
			\end{itemize} 
			
			\item Training does two things:
			\begin{itemize}
				\item Improve model performance
				\item We have loss for both 8B and 4B model -> align 4B to 8B progressively
			\end{itemize}
			
			
			 
		\end{itemize}
		
		
	\end{frame}
	
	\bibliographystyle{alpha}
	\bibliography{main}
	
\end{document}
